% Source: https://artofproblemsolving.com/community/c3h2040807_factorials
% Stable identity: aops-topic-2040807-q01
The number of trailing zeros is the exponent of \(10\), which is the smaller of the exponents of \(2\) and \(5\). In each factorial there are more factors of \(2\) than of \(5\), so we count factors of \(5\) using Legendre's formula \(v_5(n!)=\sum_{k\ge1}\left\lfloor\dfrac{n}{5^k}\right\rfloor\):
\[
\begin{array}{c|ccccccccc}
n&100&200&300&400&500&600&700&800&900\\ \hline
v_5(n!)&24&49&74&99&124&148&174&199&224
\end{array}
\]
(for example \(v_5(600!)=120+24+4=148\) and \(v_5(700!)=140+28+5+1=174\)). The total is
\[
24+49+74+99+124+148+174+199+224=1115 .
\]
So the product ends in \(\boxed{1115}\) zeros.
